\chapter{Introduction}\label{ch:intro}

\section{Context of the study}
Fast-moving consumer goods (FMCG) companies make a large number of small decisions every day. They decide how
much of each product to bake, which distributor receives which consignment, which region deserves another
sales representative and whether a promotion is worth its cost. Each decision is individually modest, but
together they determine profitability. Bakery products add a further constraint: their shelf life is short.
Over-production turns quickly into waste and under-production into empty shelves, so the value of knowing
future demand is unusually high \citep{hyndman2021}.

Dan Cake is a well-known packaged-cake brand in the Bangladeshi market. Its products range from single-serve
muffins and layer-cake slices to pound cakes, Swiss rolls, cookies and premium gift boxes, and they reach
consumers through general trade (neighbourhood shops), modern trade (supermarkets), distributors and
e-commerce. A transaction-level dataset of \NRows sales records from \DateMin to \DateMax was assembled for
this project. It is large enough that spreadsheet analysis becomes slow and error-prone, yet small enough to be
analysed in full on a laptop with open-source tools. It is therefore a realistic test bed for the framework
proposed here.

\section{Problem statement}
Routine reporting in many organisations still relies on spreadsheets prepared by hand. Such reports usually
answer only the first of three management questions: \emph{what happened?} They rarely explain \emph{why} it
happened, and they almost never estimate \emph{what is likely to happen next}. The project proposal identified
the following specific weaknesses:
\begin{itemize}
  \item large datasets are difficult to analyse manually, and manual work is slow and error-prone;
  \item historical reports do not provide forecasts, so planning is reactive;
  \item important patterns, relationships and anomalies remain hidden in the detail;
  \item managers lack an integrated, up-to-date view of key performance indicators (KPIs).
\end{itemize}
An integrated framework is therefore needed. It should process the historical data reliably, find meaningful
patterns, produce validated predictions and present them in a form that managers can act on.

\section{Aim and objectives}
The \textbf{aim} of the project is to design and develop an intelligent data analytics and machine learning
framework that transforms raw business data into actionable insights and predictive information for better
decision-making. The specific objectives are:
\begin{enumerate}[label=O\arabic*.]
  \item to collect and prepare a structured business dataset for analysis;
  \item to develop a reusable preprocessing pipeline that handles missing, duplicate, inconsistent and abnormal data;
  \item to perform exploratory data analysis (EDA) and to identify trends, relationships and patterns;
  \item to develop machine-learning models for predictive analytics and to compare several algorithms;
  \item to develop an interactive analytical dashboard;
  \item to generate business-oriented insights and recommendations from the analytical results;
  \item to evaluate the accuracy and usefulness of the framework, and to show how its outputs support management decisions.
\end{enumerate}

\section{Research questions}
\begin{enumerate}[label=RQ\arabic*.]
  \item How can data analytics techniques improve business decision-making in a sales context?
  \item Which data preprocessing techniques are most effective for preparing business transaction data?
  \item Which machine-learning algorithm gives the best predictive performance for the selected problems?
  \item What important patterns and relationships can be found in the historical sales data?
  \item How accurately can future business outcomes be predicted from historical data?
  \item How can analytical results be presented effectively to managers through an interactive dashboard?
  \item To what extent does predictive analytics improve on traditional historical reporting?
\end{enumerate}
\Cref{ch:discussion} answers each question explicitly.

\section{Scope}
The proposal allowed one business domain to be chosen. This study focuses on \textbf{sales analytics}, and
within it on four analytical problems: (i) monthly demand forecasting for each product, (ii) prediction of
large orders, (iii) segmentation of products and markets, and (iv) detection of unusual days and transactions.
The result is a working prototype, not an enterprise production system. The prototype nevertheless follows
engineering practices that would carry over to production: one configuration file, one command to run the
whole pipeline, automated tests and generated reporting.

\section{Contributions}
The main contribution is the \emph{integration} of data engineering, statistics, machine learning and business
intelligence into one reproducible pipeline. In that pipeline, every figure, table and number in the dashboard,
the website and this report is regenerated from the raw data by a single command. More specifically, the
project contributes:
\begin{itemize}
  \item a leakage-safe feature design and evaluation protocol for product-level demand forecasting, with
        time-ordered validation, multi-step (recursive) testing and comparison against naive baselines;
  \item a comparison of seven regression and seven classification algorithms, plus statistical forecasting
        models, on a real-scale business dataset;
  \item an explanation layer (price--volume decomposition, fixed-effects elasticity, permutation importance and
        SHAP values) that answers the \emph{why} question rather than reporting accuracy alone;
  \item a free, publicly deployable dashboard and project website, and an automatic insight generator that
        turns metrics into recommendations.
\end{itemize}

\section{Structure of the report}
\Cref{ch:literature} reviews the literature on business analytics, data preparation, forecasting, machine
learning and dashboards. \Cref{ch:method} sets out the research methodology. \Cref{ch:system} describes the
architecture of the framework, the dataset and the preprocessing pipeline. \Cref{ch:eda} presents the
exploratory analysis. \Cref{ch:models} reports the development and evaluation of the machine-learning models.
\Cref{ch:dashboard} describes the dashboard and the decision-support layer. \Cref{ch:discussion} discusses the
findings, answers the research questions and states the limitations. \Cref{ch:conclusion} concludes and
proposes future work.
